% !PRINTED 25
\weberpage{25}

The $a_i$ change when the $\xi$, $\eta$ are subjected to a substitution
(6). But if one forms the invariants of the form (8) (following
Algebra, Bd. I, §\,70), one obtains functions of the coefficients
which are invariant with respect to these substitutions as well, and
which therefore belong not to the individual forms $\varphi$, $\psi$,
but to the whole pencil and hence also to their intersection, the
fundamental curve. We call them \gesperrt{invariants of the fundamental curve}.
With respect to the system of forms $\varphi$, $\psi$ they are also
called \gesperrt{combinants}. The fundamental curve thus has two invariants:
\weq{9}{\begin{aligned}
A &= a_2^2 - 3\,a_1a_3 + 12\,a_0a_4,\\
B &= 27\,a_1^2a_4 + 27\,a_0a_3^2 + 2\,a_2^3 - 72\,a_0a_2a_4 - 9\,a_1a_2a_3,
\end{aligned}}
from which one derives the \gesperrt{discriminant} $D$ by the following formula:
\weq{10}{27\,D = 4\,A^3 - B^2}
With respect to the coefficients of $\varphi$ and $\psi$ the
invariants $A$, $B$, $D$ are of the degrees 8, 12, 24.

The system of forms $\varphi$, $\psi$ has \gesperrt{two simultaneous quadratic
covariants}, which are derived just as the covariant $C$ is (in Algebra I,
§\,65). The first of them, if we put
$\psi_1 = \dfrac{1}{2}\dfrac{\partial \psi}{\partial x_1}, \dots$,
is:
\weq{11}{\Phi = \begin{vmatrix}
a_{11}, & a_{12}, & a_{13}, & a_{14}, & \psi_1\\
a_{21}, & a_{22}, & a_{23}, & a_{24}, & \psi_2\\
a_{31}, & a_{32}, & a_{33}, & a_{34}, & \psi_3\\
a_{41}, & a_{42}, & a_{43}, & a_{44}, & \psi_4\\
\psi_1, & \psi_2, & \psi_3, & \psi_4, & 0
\end{vmatrix}}
and the second, $\Psi$, is obtained from it by interchanging $a_{ik}$
with $b_{ik}$ and $\psi$ with $\varphi$.

The equation $\Phi = 0$ expresses a surface of the second degree (one
that does not belong to the pencil), namely the locus of the points $x$
whose polars with respect to $\psi$ touch the surface $\varphi$.

We now assume that the discriminant $D$ is different from zero. Then
the two functions $\varphi$, $\psi$ can be transformed simultaneously
into the sum of four squares:
\weq{12}{\begin{aligned}
\varphi &= \alpha_1y_1^2 + \alpha_2y_2^2 + \alpha_3y_3^2 + \alpha_4y_4^2,\\
\psi &= \beta_1y_1^2 + \beta_2y_2^2 + \beta_3y_3^2 + \beta_4y_4^2,
\end{aligned}}
in which the $y_i$ are linear functions of the $x_i$. This is the
\gesperrt{canonical form} of the pair of functions. It then becomes:
\weq{13}{f(\xi,\eta) = (\xi\alpha_1+\eta\beta_1)(\xi\alpha_2+\eta\beta_2)(\xi\alpha_3+\eta\beta_3)(\xi\alpha_4+\eta\beta_4),}
%CONT>>
